\begin{theorem}[S-32: Two-adic integrality forces an even polarization or group degree]
\label{thm:semireg-rm-two-adic-obstruction}
Let $(X,\theta)$ be a polarized abelian
fourfold with Rosati-invariant real multiplication by
$F=\mathbf Q(\sqrt{\Delta_0})$ of type $(2,2)$, where
$\Delta_0>1$ is squarefree and
$\Delta_0\equiv2$ or $3\pmod4$.
Assume that $H_1(X,\mathbf Z_2)$ is stable under the
maximal local order
$R=\mathbf Z_2[\sqrt{\Delta_0}]$, and that every rational
divisor class has the form
$\ell(x,y)=\theta(\alpha x,y)$ for some $\alpha\in F$.
Put $D=\int_X\theta^4/24$.
The latter hypothesis is the usual minimal divisor space
$\operatorname{NS}(X)_{\mathbf Q}=F\theta$.
If $D$ is odd, then for every integral divisor class
$\ell$, every positive integer $N$, and every rational
two-adic unit $u$, the class
\[
 uC-\frac{\ell^2}{2N},\qquad
 C=\frac{\theta_1^2+\theta_2^2}{2},
 \tag{S32.1}
\]
is not two-adically integral. Consequently, if $D$ is odd,
no positive-rank perfect object has normalized character
$Nv_q$ with $Nq$ a two-adic unit.
More generally, every simple gluable fully semiregular
lift of a character $Nv_q$, $q\in\mathbf Q_{>0}$, for an effective finite coherent
target action $G$ must satisfy
\[
 2\mid D|G|.
 \tag{S32.1a}
\]
In particular, for a principal polarization there is no
nonequivariant simple gluable fully semiregular balanced-RM
perfect object, in any positive rank or for any
normalizing $B$-field.
\end{theorem}
\begin{proof}
For the integral-class assertion assume $D$ is odd.
Write $\omega=\sqrt{\Delta_0}$ and $K=\mathbf Q_2(\omega)$.
The ring $R=\mathbf Z_2[\omega]$ is a discrete valuation
ring: $x^2-\Delta_0$ is Eisenstein when $\Delta_0$ is even;
when $\Delta_0\equiv3\pmod4$, the element $\omega-1$
satisfies the Eisenstein polynomial
$x^2+2x+1-\Delta_0$.
Here we use the Eisenstein DVR criterion, for example
Sutherland, \emph{18.785 Number Theory I}, Fall 2019,
Lecture~11, Lemma~11.4 and its proof, pp.~1--2.
Put $L_2=H_1(X,\mathbf Z_2)$.
It is a finite torsion-free $R$-module, hence free of
rank four (Stacks Project, Lemma~15.22.11,
tag \texttt{0AUW}).
Set $\kappa=1/(2\omega)$.
The trace pairing $(a,b)\mapsto\operatorname{Tr}_{K/\mathbf Q_2}
(\kappa ab)$ on $R$ is perfect: its matrix in the basis
$1,\omega$ is
$\left(\begin{smallmatrix}0&1\\1&0\end{smallmatrix}\right)$.

The self-adjoint $R$-action therefore gives a unique
$R$-bilinear alternating form $\psi:L_2\times L_2\to R$
with
\[
 \theta(x,y)=\operatorname{Tr}_{K/\mathbf Q_2}
                         (\kappa\psi(x,y)).
 \tag{S32.2}
\]
For completeness, define $\psi(x,y)$ by requiring
$\theta(ax,y)=\operatorname{Tr}(\kappa a\psi(x,y))$
for all $a\in R$. Perfectness of the trace pairing gives
existence and uniqueness; self-adjointness gives
bilinearity. Also $\theta(ax,x)=0$ for every $a$, since
it equals both $\theta(x,ax)$ and $-\theta(x,ax)$;
thus $\psi(x,x)=0$ even at two.
The $\psi$-dual lattice equals the $\theta$-dual lattice:
test the trace against every $a\in R$, using $RL_2=L_2$.
Since $D$ is odd, the polarization form is unimodular
over $\mathbf Z_2$ and makes the latter lattice
$L_2$ itself. Hence $\psi$ is perfect.
Elementary symplectic elimination over the local ring
now gives an $R$-basis $e_1,f_1,e_2,f_2$ with
$\psi(e_i,f_j)=\delta_{ij}$ and the other pairings zero.
Indeed choose a primitive vector, pair it to a unit by
perfectness, rescale the second vector, and split off
its orthogonal hyperbolic plane; repeat once.

By the divisor-space hypothesis, $\ell=\theta(\alpha\cdot,\cdot)$
for $\alpha\in F$. Integrality of $\ell$ and unimodularity
of $\theta$ imply $\alpha L_2\subset L_2$.
The free $R$-basis therefore implies
$\alpha=a+b\omega$ with $a,b\in\mathbf Z_2$.
On either hyperbolic plane, the three pairings of $\ell$
corresponding to products $1,\omega,\omega^2$ are
\[
 \operatorname{Tr}(\kappa\alpha)=b,\qquad
 \operatorname{Tr}(\kappa\alpha\omega)=a,\qquad
 \operatorname{Tr}(\kappa\alpha\omega^2)=\Delta_0b.
 \tag{S32.3}
\]
Over a splitting field, $C$ is the sum of the squares
of the two embedding components of $\theta$, divided
by two. Evaluating on one ordered pair from each
hyperbolic plane consequently gives
$\operatorname{Tr}(\kappa^2\omega^j)$ when the product
of the four $R$-coefficients is $\omega^j$.
In particular, on the respective ordered quadruples
\[
 (e_1,f_1,e_2,f_2),\quad
 (\omega e_1,f_1,\omega e_2,f_2),\quad
 (\omega e_1,f_1,e_2,f_2),
\]
the values of $C$ are $1/(2\Delta_0),1/2,0$,
and those of $\ell^2/2$ are $b^2,a^2,ab$.
Thus integrality of (S32.1) would require
\[
 \frac{b^2}{N}-\frac{u}{2\Delta_0}\in\mathbf Z_2,
 \qquad
 \frac{a^2}{N}-\frac{u}{2}\in\mathbf Z_2,
 \qquad
 \frac{ab}{N}\in\mathbf Z_2.
 \tag{S32.4}
\]

Because $u$ is a two-adic unit, if $\Delta_0$ is even,
the first two congruences force
$2v_2(b)-v_2(N)=-2$ and
$2v_2(a)-v_2(N)=-1$, which give incompatible parities
for $v_2(N)$. If $\Delta_0$ is odd, both force valuation
$-1$, so $v_2(a)=v_2(b)$ and
$v_2(ab/N)=-1$, contradicting the third congruence.
This proves the integral-class assertion.

For a perfect object with normalized character
$Nv_q$, put $u=Nq$. The first two Chern equations give
$B=-\ell/N$ and
\[
 c_2(E)=\frac{N-1}{2N}\ell^2+uC
       =\frac{\ell^2}{2}
          +\left(uC-\frac{\ell^2}{2N}\right).
\]
Both $c_2(E)$ and $\ell^2/2$ are integral, contradicting
(S32.1) whenever $u$ is a two-adic unit.
Finally suppose a simple gluable fully semiregular lift
exists and $D|G|$ is odd. By S-22 its invariant Euler
pairing is
$k=4D(Nq)^2/|G|\in\{2,4\}$.
The value $k=2$ would give
$v_2((Nq)^2)=-1$, impossible for a rational square.
The value $k=4$ gives $v_2(Nq)=0$, which has just been
excluded. This proves (S32.1a) and the principal
nonequivariant corollary.
\end{proof}
